\section{A matrix Frobenius anchor at the fixed Pell parameter}
\label{sec:pell-fixed-two-transversality}

Write
\[
 (1+\sqrt2)^n=A_n+B_n\sqrt2,
 \qquad F_n(2)=B_n,\qquad L_n(2)=2A_n.
\]
The preceding Hensel section left derivative transversality at all actual
support primes as an explicit premise.  It can be discharged without using
a rank-of-apparition theorem.

Set
\[
 M=\begin{pmatrix}1&2\\1&1\end{pmatrix},\qquad
 K=\begin{pmatrix}0&2\\1&0\end{pmatrix}.
\]
Then $M=I+K$, $K^2=2I$, and the first column of $M^n$ is
$(A_n,B_n)^t$.  If $\ell=2m+1$ is prime, the intermediate binomial
coefficients are divisible by $\ell$, so for an integral matrix $R$,
\[
 M^\ell=I+K^\ell+\ell KR
       =I+2^mK+\ell KR.
\]
Taking the first column proves the following statement.

\begin{theorem}[Prime-index Frobenius anchor]
\label{thm:fixed-two-frobenius-anchor}
For every odd prime $\ell=2m+1$ there are integers $r_A,r_B$ such that
\[
 A_\ell=1+\ell r_A,
 \qquad B_\ell=2^m+\ell r_B.
\]
In particular $\ell\nmid A_\ell B_\ell$.
\end{theorem}

\begin{proof}
Only the last assertion remains.  Divisibility of the first coordinate by
$\ell$ would give $\ell\mid1$.  Divisibility of the second would give
$\ell\mid2^m$, hence $\ell=2$, contrary to oddness.
\end{proof}

The all-index differential identities specialize to
\[
 L_\ell'(2)=\ell B_\ell,
 \qquad 4F_\ell'(2)+B_\ell=\ell A_\ell.
\]

\begin{theorem}[All-support transversality at $T=2$]
\label{thm:fixed-two-all-support-transverse}
If $q\mid A_\ell$ is prime, then
$\gcd(L_\ell'(2),q)=1$.  If $p\mid B_\ell$ is prime, then
$\gcd(F_\ell'(2),p)=1$.
\end{theorem}

\begin{proof}
The coordinates are coprime.  In the first channel, $q$ divides neither
$\ell$ nor $B_\ell$, so the first differential identity proves the claim.
In the second, $p$ divides neither $\ell$ nor $A_\ell$ and is odd.  Reducing
the second identity modulo $p$ gives
$4F_\ell'(2)\equiv\ell A_\ell\not\equiv0\pmod p$.
\end{proof}

For a simple support root let $Z_1(f,t,p)$ mean that
$p\mid f(t)$, $gcd(f'(t),p)=1$, and $p^2\mid f(t)$.  The one-digit Hensel
law identifies this with zero first displacement.

\begin{theorem}[Exact fixed-parameter gate]
\label{thm:fixed-two-exact-gate}
For every odd prime $\ell$,
\[
 A_\ell B_\ell\text{ is squarefull}
 \quad\Longleftrightarrow\quad
 \begin{cases}
 Z_1(L_\ell,2,q)&(q\mid A_\ell),\\
 Z_1(F_\ell,2,p)&(p\mid B_\ell),
 \end{cases}
\]
where both quantifiers range over all rational support primes.
\end{theorem}

\begin{proof}
The preceding theorem supplies simplicity at every support prime.  The
factor $2$ in $L_\ell(2)=2A_\ell$ is a unit at every such prime, and the two
coordinates are coprime.  Thus the square divisibilities in $Z_1$ are
exactly the defining square divisibilities for squarefullness of the two
coordinates.
\end{proof}

This theorem instantiates the earlier conditional all-support statement; it
does not prove either side.  The universal exclusion of simultaneous zero
displacements is exactly the still-open prime-index squarefull exclusion.

The stronger assertion that no individual support displacement vanishes is
false.  At $\ell=7$,
\[
 A_7=239,\qquad B_7=169=13^2,
 \qquad F_7'(2)=376\equiv-1\pmod {13}.
\]
Hence $Z_1(F_7,2,13)$ holds with every premise.  This does not refute the
all-support exclusion because $239$ has exponent one.

An independent bounded replay examined all odd prime
$\ell\le20000$ and all primes $q\le10^7$ in the necessary classes
$q=\pm1\pmod{2\ell}$.  Its $3{,}091{,}963$ candidate tests found only the
same valuation-two hit $(7,13,B)$, no valuation-three hit, and no
opposite-channel repeated pair.  The 789 indices without an in-bound simple
certificate remain unresolved, as do all larger support primes.  The
finite no-hit is not an asymptotic theorem.

Lean proves Theorems~\ref{thm:fixed-two-frobenius-anchor}--
\ref{thm:fixed-two-exact-gate}.  It also proves that the surviving
zero-displacement exclusion is equivalent to squarefull exclusion, together
with the complete index-seven refutation of the stronger local assertion.
The standard abc conjecture and the fixed-Pell squarefull exclusion both
remain open.
